\documentclass[BCOR=6mm,headmarks=true,automark,font=palatino,
  headline=oneliner,parskip=half]{tumbook}


\usepackage{lipsum}
\usepackage{booktabs}
\usepackage{esint}
\usepackage{cleveref}


\title{Computational Methods in Nanoelectronics}
\subtitle{The simulation of everything}
\edition{Summer Term 2021}

\dedication{To Franz X. Gabelsberger, inventor of the street named after him.}

\renewcommand{\theTUMCoverImage}{%
  \includegraphics[width=\TUMCoverImageWidth,
    height=\TUMCoverImageHeight]{example-image}
}


\begin{document}

\frontmatter
\maketitle

\chapter{Preface}
 {\color{TUMOrange}This is a version of the \emph{tumbook} class using
  \textbf{headmarks} and the \textbf{palatino} font.
  It contains a \emph{coverpage} with title, subtitle, edition, and author.
  A custom cover image is configured. The maximum size of the cover image is
  half of the page. The coverpage uses the TUM  titlepage layout.
  A \textbf{bindingoffset of 6mm} is specified for the Document.\par
  The following \emph{titlepage} contains the title information. No publisher
  or titleback information is given. The next page contains the
  \emph{dedication}.\par
  After the \emph{preface} the \emph{table of contents},
  \emph{list of figures}, and \emph{list of tables} follow.
  The page numbering of the entire \emph{frontmatter} is printed as
  lowercase Roman numerals.\par
  The \emph{mainmatter} contains an example document with multiple chapters,
  sections, and subsections. The page numbering is printed in Arabic numbers
  always at the outer margin of the page. \textbf{palatino} is chosen as the
  main text font. This book uses \textbf{headmarks} printing the chapter name
  on left pages and the section name on right pages.
 }

\lipsum[1-4]
\cleardoublepage

\tableofcontents
\listoffigures
\listoftables

\mainmatter
\chapter{The beginning of the book}
\section{The first section}
\lipsum[1-2]\cite{jirauschek2014}

\section{Another section, this time with a long title}
\lipsum[3-4]
\begin{figure*}[tbp]
  \centering
  \includegraphics[width=.8\textwidth, height=.3\textheight]{example-image-a}
  \caption{This is an example image that should be printed
    at the top of a page.}
  \label{fig:one}
\end{figure*}
\lipsum[5-7]

\subsection{Sections can have subsections}
\lipsum[8-10]

\subsection{Another subsection}
\lipsum[11]
\begin{equation}
  \label{eq:one}
  a^2 + b^2 = c^2
\end{equation}
\lipsum[12]

\section{Some nice theory}
\lipsum[13]\cite{feynman1958theory}

\subsection{A particular use-case}
\lipsum[14]
\begin{table}
  \centering
  \caption{This is a table caption at the top of a table.}
  \label{tab:one}
  \begin{tabular}{lcr}
    \toprule
    Header 1 & Header 2 & Header 3 \\
    \midrule
    Data 1   & 3.5      & 4.3      \\
    Data 2   & 3.5      & 4.3      \\
    Data 3   & 3.5      & 4.3      \\
    Data 4   & 3.5      & 4.3      \\
    Data 5   & 3.5      & 4.3      \\
    \bottomrule
  \end{tabular}
\end{table}
\lipsum[15-18]

\subsection{Problems with the theory}
\lipsum[19-25]

\section{A better theory}
\lipsum[26-28]\cite{bardeen1973four}
\begin{figure*}[tbp]
  \begin{subfigure}[t]{.4\textwidth}
    \centering
    \includegraphics[width=0.8\textwidth]{example-image-b}
    \subcaption{Subfigure left, with its own caption.}
    \label{fig:two_a}
  \end{subfigure}
  \hfill
  \begin{subfigure}[t]{.5\textwidth}
    \centering
    \includegraphics[width=0.8\textwidth]{example-image-c}
    \subcaption{Subfigure right, with its own caption
      \cite{einstein1905elektrodynamik}.}
    \label{fig:two_b}
  \end{subfigure}
  \centering
  \caption{Another caption for the entire figure. The figure consists of
    two subfigures.}
  \label{fig:two}
\end{figure*}

\subsection{A long subsection title that spans multiple lines is not a good
  idea but it is possible nevertheless}
\lipsum[29-30]
\begin{align}
  \nabla\cdot\vec{D}  & = \rho                                         \\
  \nabla\cdot\vec{B}  & = 0                                            \\
  \nabla\times\vec{E} & = -\frac{\partial\vec{B}}{\partial t}          \\
  \nabla\times\vec{H} & = \vec{J} + \frac{\partial\vec{D}}{\partial t}
\end{align}
\lipsum[31]
\begin{gather}
  \oiint_{\partial\Omega}\vec{D}\times\mathrm{d}\vec{S}
  = \iiint_\Omega\rho\mathrm{d}V                            \\
  \oiint_{\partial\Omega}\vec{B}\times\mathrm{d}\vec{S} = 0 \\
  \oint_{\partial\Sigma}\vec{E}\mathrm{d}l
  = -\frac{d}{dt}\iint_\Sigma\vec{B}\mathrm{d}\vec{S}       \\
  \oint_{\partial\Sigma}\vec{H}\mathrm{d}l
  = \iint_\Sigma\vec{J}\mathrm{d}\vec{S}
  +\frac{d}{dt}\iint_\Sigma\vec{D}\mathrm{d}\vec{S}
\end{gather}
\lipsum[32-35]

\subsection{The last subsection of this chapter}
\lipsum[36-40]


\chapter{Mathematical Foundations}
\section{Linear Algebra}
\label{sec:linalg}
\lipsum[1-2]

\subsection{Vector spaces}
\lipsum[3-4]
\begin{tcolorbox}[TUMBox=TUMGreen, title=Definition, breakable]
  Sei $K$ ein Körper. Eine Menge $V$ zusammen mit zwei Verknüpfungen
  \begin{equation}
    +: V \times V \rightarrow V, (a,b) \rightarrow a+b \quad\text{und}\quad
    \cdot: K \times V \rightarrow V, (\lambda,b) \rightarrow \lambda\cdot b
  \end{equation}
  heißen $\boldmath{K}$\textbf{-Vectorraum} wenn gilt:\par

  \begin{description}[itemsep=0pt]
    \item[(V1)] Die Verknüpfung $+$ ist assoziativ, das heißt:
      \begin{equation}
        (x+y)+c = x+(y+z) \quad\forall x,y,z \in V.
      \end{equation}
    \item[(V2)] Die Verknüpfung $+$ ist kommutativ, das heißt:
      \begin{equation}
        x+y=y+x \quad\forall x,y \in V.
      \end{equation}
    \item[(V3)] Es gibt ein neutrales Element bezüglich $+$, genannt $0_V$
      (Nullvector), das heißt:
      \begin{equation}
        0_V+v=v+0_V=v \quad\forall v \in V.
      \end{equation}
    \item[(V4)] Jedes Element $v \in V$ besitz ein inverses Element bezüglich
      $+$, geschrieben $-v$, das heißt:
      \begin{equation}
        v+(-v)=-v+v=0 \quad\forall v \in V.
      \end{equation}
    \item[(V5)] Die Verknüpfung $\cdot$ ist assoziativ, das heißt:
      \begin{equation}
        \lambda\cdot(\mu\cdot v) = (\lambda\mu)\cdot v \quad\forall\lambda,\mu
        \in K \forall v \in V.
      \end{equation}
    \item[(V6)] Es gibt ein neutrales Element bezüglich $\cdot$, genannt das
      Einselement $1$ des Körpers $K$, das heipt:
      \begin{equation}
        1\cdot v = v\cdot 1 \quad\forall v \in V.
      \end{equation}
    \item[(V7)] Es gilt das Distributivgesetz:
      \begin{equation}
        \lambda\cdot(v+w) = \lambda\cdot v + \lambda\cdot w \quad\forall
        \lambda\in K, \forall v \in V.
      \end{equation}
  \end{description}
  Die Elemente $v \in V$nennen wir nun \textbf{Vektoren}. Die Verknüpfung
  $+$ heißt \textbf{Vektoraddition} und $\cdot$ heißt
  \textbf{Skalarmultiplikation}.
\end{tcolorbox}
\lipsum[5-6]
\begin{TUMBoxFill}[TUMOrange]{Alert}
  This is an improtant fact to remember!
\end{TUMBoxFill}
\lipsum[7]

\subsection{Matrices and linear transformations}
\lipsum[8-12]

\subsection{Eigenvectors and Eigenvalues}
\lipsum[13-20]

\section{Analysis}
\label{sec:analysis}
\lipsum[21-22]

\subsection{Differentials}
\lipsum[23-25]
\begin{tcolorbox}[TUMBox, title=Of utmost importance]
  \lipsum[1]
  \begin{equation}
    \frac{d}{dx}x^n = nx^{n-1}
  \end{equation}
  \lipsum[2]
\end{tcolorbox}
\lipsum[26-29]

\subsection{Integrals}
\lipsum[30-35]

\section*{Summary}
\label{sec:summary}
\lipsum[36]
\begin{figure}[tbp]
  \centering
  \includegraphics[width=.4\textwidth]{example-image}
  \caption{Another very important figure.}
\end{figure}
\lipsum[37-40]

\section*{Exercises}
\label{sec:exercises}
\lipsum[41-43]


\appendix
\chapter{Appendix}
\section{Supporting Data}
\lipsum[1]
\begin{table}[htb]
  \centering
  \caption{This table contains some supporting data related to
    \cref{sec:linalg}.}
  \label{tab:supp}
  \begin{tabular}{llll}
    \toprule
    A      & B   & C   & D    \\
    \midrule
    Data 1 & 3.5 & 4.3 & 1.12 \\
    Data 2 & 3.5 & 4.3 & 1.12 \\
    Data 3 & 3.5 & 4.3 & 1.12 \\
    Data 4 & 3.5 & 4.3 & 1.12 \\
    Data 5 & 3.5 & 4.3 & 1.12 \\
    \bottomrule
  \end{tabular}
\end{table}

\section{Trivial Proof}
\lipsum[2-6]


\backmatter
\bibliographystyle{IEEEtran}
\bibliography{content/literature}

\end{document}
